% Fate's Edge – The Pact-Whisperer’s Guide
% A standalone supplement for summoners
% Compile with XeLaTeX or LuaLaTeX
\documentclass[10pt,a5paper]{book}
\usepackage[top=1.5cm,bottom=1.5cm,left=1.2cm,right=1.2cm]{geometry}
\usepackage{fontspec}
\setmainfont{TeX Gyre Pagella}
\usepackage{titlesec}
\usepackage{tcolorbox}
\tcbuselibrary{skins,breakable}
\usepackage{multicol}
\usepackage{hyperref}
\hypersetup{colorlinks=true,linkcolor=blue}
\usepackage{array}
\usepackage{longtable}
\usepackage{epigraph}
\setlength{\epigraphwidth}{0.7\textwidth}
\renewcommand{\epigraphflush}{flushright}
\usepackage{pgfornament}
\usepackage{lettrine}

% Chapter styling
\titleformat{\chapter}[display]
{\Huge\bfseries\scshape}
{\flushright\pgfornament[width=0.2\textwidth]{88}\\\vspace{-1cm}}
{0pt}
{\hfill}
\titlespacing*{\chapter}{0pt}{-20pt}{40pt}

% Section styling
\titleformat{\section}
{\Large\bfseries}
{\thesection}{0.5em}{}
\titlespacing*{\section}{0pt}{12pt}{6pt}

% Subsection
\titleformat{\subsection}
{\normalsize\bfseries\itshape}
{\thesubsection}{0.5em}{}
\titlespacing*{\subsection}{0pt}{8pt}{4pt}

% Custom environments
\newtcolorbox{patronbox}[2][]{
    colback=gray!5,
    colframe=gray!70,
    fonttitle=\bfseries,
    title=⛧ #2 ⛧,
    enhanced,
    attach boxed title to top left={yshift=-2mm,xshift=3mm},
    boxed title style={colback=gray!20},
    #1
}

\newtcolorbox{spiritcard}[2][]{
    breakable,
    colback=blue!5,
    colframe=blue!70,
    fonttitle=\bfseries\large,
    title=#2,
    enhanced,
    attach boxed title to top left={yshift=-3mm,xshift=3mm},
    boxed title style={colback=blue!20},
    #1
}

\newtcolorbox{ritualnote}{
    colback=yellow!5,
    colframe=orange!70,
    before skip=0.5cm,
    after skip=0.5cm
}

\newcommand{\ritualtag}[1]{\textsc{#1}}

\begin{document}

\frontmatter
\thispagestyle{empty}
\vspace*{2cm}
\begin{center}
{\Huge\scshape The Pact-Whisperer’s Guide}\\[0.5cm]
{\Large Binding Spirits to Will}\\[1cm]
\pgfornament[width=0.6\textwidth]{88}\\[1cm]
{\itshape “I do not command the dead. I ask. And sometimes – when the price is right – they answer.”}
\end{center}
\vfill
\clearpage

\tableofcontents
\clearpage

\mainmatter

\chapter{Introduction – The Gray Wanderer’s Confession}
\epigraph{“I have called spirits from the deep stone, from the salt foam, from the ashes of a pyre. Some came as friends. Most came as bargains. A few came as mistakes I still dream about.”}{\textit{— The Gray Wanderer}}

This is not a grimoire of safe magic. Summoning is the art of opening doors and hoping you close them before something follows you through. Every culture has its own way – contracts, offerings, drums, salt circles, cursed coins – but the heart is the same: \textit{you need something the dead or the never‑born can give}.

In these pages, I have gathered traditions from across the Broken Marches to the Zakov reefs. You will find spirits that mend your boots and spirits that drown your enemies. You will learn how a Ykrul bone‑singer honours a wolf‑ancestor, and how a Zakov fathom‑caller pays a drowned mutineer in blood coins. You will also learn when to walk away.

Because the spirits remember. And they have long memories.

\clearpage

\chapter{The Lore – How Cultures Summon}
\epigraph{“Every people has a door. The trick is knowing which key to use – and whether you are willing to pay the toll.”}{\textit{— The Gray Wanderer}}

The following traditions are not exclusive – a Linn tide‑rider might learn a dwarf’s runic contract, and a halfling hearth‑caller might borrow a Ykrul drum. But each culture’s native method reflects its relationship with the unseen. Pay attention to the \textbf{Patrons} mentioned in passing; they often lurk behind the rites, watching.

\section{Aeler (Dwarves) – The Stone-Whisperers}
\epigraph{“The mountain does not speak to those who shout. It speaks to those who carve patiently, and listen for the echo.”}{\textit{— Runekeeper Durin}}

\lettrine{D}{warves} do not summon spirits from elsewhere. They \textit{awaken} spirits that already sleep in the stone – the memory of a mountain, the grudge of a collapsed tunnel, the patience of a deep vein. These are not Outsiders in the usual sense; they are \textit{indigenous} to the rock, and have been there since the first hammer struck flint.

\subsection*{The Rite of the Carved Contract}
A dwarf summoner spends an hour carving a contract onto a slate tablet. The contract specifies:
\begin{itemize}
    \item The spirit’s name (as known to the mountain, often in Old Aeler runes)
    \item The task (e.g., “guard this door for three days” or “report when strangers approach”)
    \item The payment (e.g., “a vein of copper to be left untouched” or “a song sung in the deep hall”)
\end{itemize}
The summoner then taps the tablet with a stone hammer. If the mountain accepts, the tablet warms and the spirit emerges from the nearest stone surface – a wall, a floor, a boulder.

\subsection*{Patrons and the Deep Stone}
The Aeler revere the \textbf{Nidhoggr} (World‑Worm) as the dreaming ancestor of all stone. They also whisper to the \textbf{Sealed Gate} when they need a spirit contained, and to \textbf{Mykkiel} when the contract must be witnessed. A dwarf summoner who ignores these Patrons finds that stone grows silent.

\subsection*{Cultural Taboos}
\begin{itemize}
    \item Never summon a spirit whose name you have not heard from three independent sources. The mountain does not forgive misnaming.
    \item Destroy the contract tablet after the spirit departs, or the spirit may claim you as its next task.
    \item Do not summon in a mine that has recently claimed a life. The spirit will be that miner, and they will be furious.
\end{itemize}

\subsection*{Sample Spirit – Vein-Spirit (Lesser, Cap 1)}
\begin{spiritcard}{Vein-Spirit (Lesser)}
A small, metallic serpent that slithers through rock.
\begin{itemize}
    \item \textbf{Cap:} 1 (Leash = 1 + Spirit of summoner)
    \item \textbf{Harm:} 2 segments
    \item \textbf{Pools:} Physical 3d (Body 2 + Cap 1); Magical 2d
    \item \textbf{Abilities:} Locate a specific mineral within 100 metres (grant +2 dice to Mining or Prospecting), seal a small crack in stone (treat as \textit{Mend} tag for one scene)
    \item \textbf{Instinct:} Coils around the richest ore and refuses to move.
    \item \textbf{Ban:} A copper coin hammered flat. Press it to the stone, and the spirit will retreat.
\end{itemize}
\end{spiritcard}

\subsection*{Story Hook}
\epigraph{“I carved a contract for a Vein‑Spirit to find me a new iron lode. It found one – under the old cemetery. The town council is not happy. The spirit, however, is very happy.”}{\textit{— Aeler Prospector}}

\clearpage

\section{Aelaerem (Halflings) – The Hearth-Callers}
\epigraph{“A full belly and a warm fire are better than any binding circle. Offer bread, and the spirits will offer forgiveness – usually.”}{\textit{— Aelaerem Hearth-Mother}}

\lettrine{H}{alflings} summon small things – brownies, hearth‑spirits, field‑wights – through offerings of bread and milk. They do not bind; they \textit{befriend}. A halfling summoner’s Leash is replaced by a \textit{Friendship} clock. When it fills, the spirit does not attack – it simply leaves, slightly disappointed, and may refuse to answer again for a year and a day.

\subsection*{The Rite of the Warm Saucer}
Place a saucer of warm milk and a crust of fresh bread on the hearthstone. Knock three times and say, “Little friend, little neighbour, I have a small task and a small thanks.” If a spirit accepts, it will appear within ten minutes – usually as a shadow, a flicker of light, or a rustle in the thatch.

\subsection*{Patrons of the Threshold}
The Aelaerem honour \textbf{Inaea} (Angel of the Spider) as the weaver of small bonds, and \textbf{Morag} (Weaver of Hidden Costs) as the one who remembers every promise. \textbf{Palinode} is invoked when the spirit’s task involves song or storytelling. But the true patron of hearth‑callers is \textbf{the Traveler} – for what is a hearth if not a waystation?

\subsection*{Cultural Taboos}
\begin{itemize}
    \item Never summon a hearth‑spirit if your own hearth has been cold for more than a day. They will think you are a bad neighbour.
    \item Do not offer soured milk. The spirit may sour your butter for a month.
    \item If a field‑wight accepts your offering but you fail to complete its task, you must leave a larger offering on your doorstep every night for a year.
\end{itemize}

\subsection*{Sample Spirit – Brownie (Lesser, Cap 1)}
\begin{spiritcard}{Brownie (Lesser)}
A small, gnarled figure that stands six inches tall.
\begin{itemize}
    \item \textbf{Cap:} 1 (Friendship clock = 4 segments; when filled, spirit leaves kindly)
    \item \textbf{Harm:} 1 segment (brownies flee rather than fight)
    \item \textbf{Pools:} Physical 2d, Magical 3d (subtle magic)
    \item \textbf{Abilities:} Clean a room in an hour, untie any knot (automatic), hide a small object so that no one can find it for a day
    \item \textbf{Instinct:} Mends socks and polishes boots, even when not asked.
    \item \textbf{Ban:} A whistle at midnight – it will flee, thinking the Wild Hunt is near.
\end{itemize}
\end{spiritcard}

\subsection*{Story Hook}
\epigraph{“Our brownie hid the tax collector’s quill. The collector wrote ‘paid’ in blood and left. The brownie now expects a new hat every spring.”}{\textit{— Aelaerem Innkeeper}}

\clearpage

\section{Ykrul (Orcs) – The Bone-Singers}
\epigraph{“The drum is not wood and hide. It is the heart of the ancestor, still beating. You do not play it – you listen to it.”}{\textit{— Bone-Singer Khelt}}

\lettrine{Y}{krul} summon the spirits of honoured ancestors or totem beasts. They sing a name over a drum made from that creature’s hide. The spirit arrives angry, proud, and requires a show of respect before it will serve.

\subsection*{The Rite of the Drummed Name}
The summoner drums a rhythm – the heartbeat of the creature – while singing the spirit’s name. The drum must be made from the same species as the spirit (wolf hide for a wolf ancestor, bear hide for a bear totem). After the third verse, the spirit appears, usually as a translucent version of the living animal or a humanoid ancestor wearing its fur.

\subsection*{Patrons of the Steppe}
Ykrul bone‑singers answer to \textbf{Aveh} (the Rider Behind the Storm) when summoning wind‑spirits, and to \textbf{Grimmir} (the Wild Speaker) for totem beasts. But the greatest patron is \textbf{Varnek Karn} – the Death’s Negotiator. Every ancestor is a dead thing, and Varnek Karn brokers the terms.

\subsection*{Cultural Taboos}
\begin{itemize}
    \item Never summon an ancestor whose name you have not earned the right to speak. The clan keeps a list.
    \item Do not use a drum made from an animal you did not kill yourself. The spirit will know.
    \item If the spirit asks, “Who calls me?” you must answer with your full lineage, back seven generations. Missing a name weakens the binding.
\end{itemize}

\subsection*{Sample Spirit – Wolf Ancestor (Greater, Cap 3)}
\begin{spiritcard}{Wolf Ancestor (Greater)}
A towering grey wolf with silver eyes.
\begin{itemize}
    \item \textbf{Cap:} 3 (Leash = 3 + Spirit)
    \item \textbf{Harm:} 4 segments
    \item \textbf{Pools:} Physical 6d, Magical 3d
    \item \textbf{Abilities:} Track any scent for up to a day (automatic success on Survival), grant +2 dice to all allies’ Perception for a scene, deal Harm 2 (bite) with a successful attack
    \item \textbf{Instinct:} Hunts the largest target in sight.
    \item \textbf{Ban:} A howl sung off‑key – the spirit will cover its ears and retreat.
\end{itemize}
\end{spiritcard}

\subsection*{Story Hook}
\epigraph{“I called my great‑grandmother’s wolf. She appeared, teeth bared. ‘You are thin,’ she said. ‘You have not hunted in a moon.’ I bowed. ‘I hunt a traitor now. Help me.’ She found the trail in seconds. Then she demanded a haunch of the kill. I gave it. She vanished, satisfied. But now she visits every full moon, expecting more.”}{\textit{— Bone-Singer Khelt}}

\clearpage

\section{Linn (Northmen) – The Tide-Riders}
\epigraph{“The sea has more ghosts than the land. Every wave carries a whisper. Every reef is a graveyard.”}{\textit{— Linn Tide-Rider}}

\lettrine{L}{inn} summon sea spirits – selkies, wave‑walkers, drowned crewmates. They call them at the water’s edge, using a knotted rope as a focus. The spirit arrives wet and whispering, speaking in the language of tides.

\subsection*{The Rite of the Nine Knots}
Find a place where salt water meets land – a beach, a tidal pool, a river mouth. Tie nine knots in a hemp rope, each knot corresponding to one of the nine waves of a storm swell. As you tie each knot, speak the spirit’s name. On the ninth knot, cast the rope into the water. The spirit will rise from the foam.

\subsection*{Patrons of Salt and Sorrow}
The Linn revere \textbf{Raéyn} (Mistress of the Sea) as the mother of all tide‑spirits, but they whisper to \textbf{Khemesh} (the Abyssal Maw) when they need a drowning’s power. \textbf{Palinode} is invoked for sea shanties that call selkies, and \textbf{the Traveler} for safe passage across the whale‑road.

\subsection*{Cultural Taboos}
\begin{itemize}
    \item Never summon a sea spirit during a dead calm – they cannot reach you without waves.
    \item Do not use a rope made from land fibres (cotton, linen) – only hemp or raw wool. Synthetics insult them.
    \item If the spirit asks for a drink, pour a cup of fresh water into the sea. Do not use salt water – they have enough of that.
\end{itemize}

\subsection*{Sample Spirit – Selkie (Greater, Cap 3)}
\begin{spiritcard}{Selkie (Greater)}
A seal that can shed its skin to become a beautiful human.
\begin{itemize}
    \item \textbf{Cap:} 3 (Leash = 3 + Spirit)
    \item \textbf{Harm:} 4 segments
    \item \textbf{Pools (seal form):} Physical 6d (swimming), Magical 4d; (human form) Physical 4d, Magical 5d
    \item \textbf{Abilities (seal):} Swim faster than any ship (automatically win a chase at sea), hold breath for an hour. \textbf{(human):} Sing a song of forgetting (target forgets last 5 minutes, Resolve DV 4)
    \item \textbf{Instinct:} Wants to find its lost sealskin and return to the sea.
    \item \textbf{Ban:} A circle of salt – it cannot cross, as salt reminds it of shore.
\end{itemize}
\end{spiritcard}

\subsection*{Story Hook}
\epigraph{“I called a selkie to guide us through the reef. She came as a woman, dripping wet. ‘What will you give?’ ‘A silver brooch.’ She laughed. ‘That will buy my voice for an hour.’ She led us through a channel I had never seen. Then she took her sealskin from a pouch, changed, and swam away. My brooch was gone from my hand – and my left ear has whispered ever since.”}{\textit{— Linn Tide-Rider}}

\clearpage

\section{Zakov (Pirates) – The Salt-Binders}
\epigraph{“A drowned man’s coin is heavier than gold. It pulls you down – but first, it pulls secrets up.”}{\textit{— Fathom-Caller Sable}}

\lettrine{Z}{akov} summoners (often called \textit{Fathom-Callers}) use cursed coins and drowned names. They specialise in binding the vengeful dead – pirates who were betrayed, captains who were mutinied. These spirits are dangerous, powerful, and always hungry for revenge.

\subsection*{The Rite of the Swallowed Coin}
Take a coin that has been used to pay a murderer (a “blood coin”). Place it on your tongue and swallow it (mark 1 Fatigue). Then whisper the name of a drowned person three times. The spirit will crawl out of the nearest shadow, wet and hissing.

\subsection*{Patrons of the Deep Debt}
Zakov salt‑binders serve \textbf{Malachai} (the Cruel Messenger) – every cursed coin is a tiny corruption. They also bargain with \textbf{Maelstraeus} (the Infernal Bargainer), for is not a coin a contract? Few invoke \textbf{the Witness} – the truth of their bargains would damn them.

\subsection*{Cultural Taboos}
\begin{itemize}
    \item Never summon a spirit you cannot offer a specific revenge target. The spirit will default to \textit{you}.
    \item Do not swallow the coin if it has been used to pay a priest – that coin is blessed and will choke you (Harm 1).
    \item If the spirit asks, “Whose blood?” you must name someone. Lying causes the spirit to attack immediately.
\end{itemize}

\subsection*{Sample Spirit – Drowned Mutineer (Lesser, Cap 1)}
\begin{spiritcard}{Drowned Mutineer (Lesser)}
A skeletal figure in rotting sailor’s clothes.
\begin{itemize}
    \item \textbf{Cap:} 1 (Leash = 1 + Spirit)
    \item \textbf{Harm:} 2 segments
    \item \textbf{Pools:} Physical 3d (clawing), Magical 2d (curses)
    \item \textbf{Abilities:} Curse a rope to snap at a critical moment (treat as 1 SB), whisper a secret into a sleeping target’s ear (target wakes with a new Lead for an investigation)
    \item \textbf{Instinct:} Tries to untie anything tied – knots, ropes, shoelaces, marriages.
    \item \textbf{Ban:} A bell rung three times (the signal for “abandon ship”).
\end{itemize}
\end{spiritcard}

\subsection*{Story Hook}
\epigraph{“I called a captain who had been thrown overboard by his own crew. He appeared, eye sockets full of barnacles. ‘You have the coin?’ I nodded. ‘Then listen. My first mate is now a harbourmaster in Silkstrand. Bring me his gold – not his blood – and I will guide your ship through the Straits for a year.’ I did not make that bargain. I burned the coin instead. The captain screamed, but he could not reach me through the fire. Now I hear him scratching at my hull every night.”}{\textit{— Fathom-Caller Sable}}

\clearpage

\section{Other Traditions – Brief Notes}
\epigraph{“There are as many ways to call a spirit as there are stars. These are only the ones I have seen with my own eyes – and regretted.”}{\textit{— The Gray Wanderer}}

\begin{itemize}
    \item \textbf{Valewood (Lethai‑ar):} \textit{Root-Weavers}. Use living wood and sap offerings to call forest shades and beast‑kin. Their Patron is \textbf{Grimmir} (the Wild Speaker). Ban: iron tools in the grove.
    \item \textbf{Mistlands:} \textit{Bell-Wakes}. Ring a bell tuned to the spirit’s former note; the spirit rises as a mist‑wraith. Patrons: \textbf{the Sealed Gate} (to keep it contained) and \textbf{Varnek Karn} (to speak with the dead). Ban: a lantern lit with wick‑silk from a living person’s hair.
    \item \textbf{Acasia (Broken Marches):} \textit{Curse-Summoners}. Call the ghosts of famine victims using a handful of rotten grain. Patron: \textbf{the Carrion‑King} (decay and renewal). Ban: a loaf of fresh bread offered before the summoning – the spirit cannot abide plenty.
\end{itemize}

\clearpage

\chapter{The Spirit Compendium}
\epigraph{“I have met a spirit of the forge who only spoke in hammer strikes, and a spirit of the hearth who sang lullabies. They were both more honest than most people I know.”}{\textit{— The Gray Wanderer}}

This compendium expands the sample spirits from the cultures above and adds general spirits usable by any summoner. Each entry includes the spirit’s typical attitude toward Patrons – useful for roleplaying.

\section{Elemental Spirits}

\begin{longtable}{p{3cm} p{2.5cm} p{8cm}}
\toprule
Name & Cap & Description \\
\midrule
Emberling & Lesser (1) & Small fire spirit. Can ignite unattended objects, provide warmth, or flare as a distraction (+1 die to Deception). \textbf{Patron:} Palinode (performance of flame). \textbf{Instinct:} burns paper and dry grass. \textbf{Ban:} a handful of wet ash. \\
Cinder-Wraith & Greater (3) & Ash‑and‑flame humanoid. Can shoot a burning bolt (Harm 1, Near), leave a trail of smoke (concealment). \textbf{Patron:} Malachai (cursed fire). \textbf{Instinct:} spreads fire to the nearest structure. \textbf{Ban:} a bell rung three times. \\
Magma-Heart & Elder (5) & Core of liquid stone. Can melt passages, cause tremors (Area knockdown), radiate lethal heat (Harm 2 to Close each round). \textbf{Patron:} Nidhoggr (stone’s wrath). \textbf{Instinct:} seeks the lowest point to pool. \textbf{Ban:} a bucket of sea water (cools, does not banish). \\
\bottomrule
\end{longtable}

\section{Ancestral Spirits}

\begin{longtable}{p{3cm} p{2.5cm} p{8cm}}
\toprule
Name & Cap & Description \\
\midrule
Shade of the Forge & Lesser (1) & Ancestral dwarf smith. Repairs a metal object (remove Compromised), or grants +1 die to a Craft roll. \textbf{Patron:} Mykkiel (contract of craft). \textbf{Instinct:} demands fuel (coal, wood). \textbf{Ban:} a flawless piece of unworked iron. \\
Battle-Howl Echo & Greater (3) & Roaring ancestor spirit. Inspires allies (+1 die to Fear resist), frightens enemies (–1 die to attacks). \textbf{Patron:} Oath of Flame \& Light (battle‑vows). \textbf{Instinct:} charges the largest enemy. \textbf{Ban:} a white flag or verbal surrender. \\
Stone-Father & Elder (5) & Mountain’s memory. Reshapes stone, speaks with stone‑dead, crushing aura (enemies half speed). \textbf{Patron:} Nidhoggr. \textbf{Instinct:} settles into stable position and refuses to move. \textbf{Ban:} a seismic shock (harm 3 to spirit, may collapse cave). \\
\bottomrule
\end{longtable}

\section{Sea Spirits}

\begin{longtable}{p{3cm} p{2.5cm} p{8cm}}
\toprule
Name & Cap & Description \\
\midrule
Skittering Foam & Lesser (1) & Tiny wave‑creature. Douses flames, cleans an object, carries a small object (under 1 kg) across water. \textbf{Patron:} Raéyn. \textbf{Instinct:} seeks the lowest point. \textbf{Ban:} a handful of dry sand. \\
Drowned Helmsman & Greater (3) & Ghost of a lost pilot. Guides ship (+2 dice navigation), creates fog (Near concealment). \textbf{Patron:} Khemesh (inevitable tide). \textbf{Instinct:} steers toward former home port. \textbf{Ban:} a coin placed on its eyes. \\
Kraken’s Spite & Elder (5) & Fragment of Khemesh’s anger. Grapples ship, crushes hull, summons a tentacle (Cap 1). \textbf{Patron:} Khemesh. \textbf{Instinct:} drags everything toward deep water. \textbf{Ban:} a prayer to Raéyn with genuine sorrow. \\
\bottomrule
\end{longtable}

\clearpage

\chapter{The Summoner’s Talents}
\epigraph{“A talent is not a gift. It is a scar that healed in a useful shape.”}{\textit{— The Gray Wanderer}}

These talents expand the summoning rules. Some are available to any summoner; others require specific Paths.

\section{Minor Talents (2–5 XP)}

\begin{tcolorbox}[colback=green!5, colframe=green!70]
\textbf{Boon Finesse} (2 XP, any summoner) \\
Once per round, spend 1 Boon to clear 1 tick from your active spirit’s Leash. Cannot be used after the Leash fills.
\end{tcolorbox}

\begin{tcolorbox}[colback=green!5, colframe=green!70]
\textbf{Honourable Departure} (3 XP, any summoner) \\
When you voluntarily dismiss a spirit (action, before Leash fills), gain +1 Boon and the next summoning of the same spirit reduces Call DV by 1 (it remembers courtesy).
\end{tcolorbox}

\begin{tcolorbox}[colback=green!5, colframe=green!70]
\textbf{Sympathetic Focus} (4 XP, any summoner) \\
Use a physical token tied to the spirit (scale, coal, bone) as a focus. Your Call action is one step faster (Some Time → Moment, etc.).
\end{tcolorbox}

\section{Major Talents (6–10 XP)}

\begin{tcolorbox}[colback=blue!5, colframe=blue!70]
\textbf{Spirit Link} (10 XP, any summoner) \\
Your spirit acts on your turn, not its own initiative. Commands are free actions. The spirit moves and acts immediately. Reduce Leash ticking for natural behaviours by 1 per round.
\end{tcolorbox}

\begin{tcolorbox}[colback=blue!5, colframe=blue!70]
\textbf{Greater Binding} (8 XP, Spirit 3+) \\
May summon and maintain a Greater or Elder spirit without increasing Call DV penalty for multiple spirits. However, each spirit beyond the first adds +1 DV to all summoning rolls.
\end{tcolorbox}

\begin{tcolorbox}[colback=blue!5, colframe=blue!70]
\textbf{Soul Anchor} (8 XP, Codex or Symbol required) \\
Your Codex or Symbol anchors one spirit. That spirit does not depart at downtime; you may Call it instantly (no action, no DV) once per session. The anchor must touch you.
\end{tcolorbox}

\section{Prestige Talents (12+ XP)}

\begin{tcolorbox}[colback=red!5, colframe=red!70]
\textbf{Pact-Broker} (14 XP, Tier III, Spirit 4+) \\
Negotiate a \textit{Pact} with a spirit: a specific task, a time limit, a price. While the spirit works toward the pact, it does not tick Leash. Breach of pact (either side) causes Harm 3 and immediate unbinding. Dangerous – use only when trust is high.
\end{tcolorbox}

\begin{tcolorbox}[colback=red!5, colframe=red!70]
\textbf{Throne of the Bound} (16 XP, Tier IV, Greater Binding, Spirit 5+) \\
Control up to three spirits simultaneously (one may be Elder). Call action replaced by \textit{Chorus} (DV = highest Cap + number of spirits). All share one Leash track (Cap = Spirit+Presence). When it fills, all act to their natures at once – chaos assured.
\end{tcolorbox}

\clearpage

\chapter{Summoning in Play – GM’s Guide}
\epigraph{“The spirit is not a tool. It is a person who happens to be dead, or never born, or made of fire. Treat it badly, and it will remember.”}{\textit{— The Gray Wanderer}}

\section{When to Tick the Leash – Examples}
\begin{itemize}
    \item \textbf{Tick 1:} The spirit uses a powerful ability (e.g., Emberling’s ignition).
    \item \textbf{Tick 1:} The summoner takes a significant action while the spirit is active (attacking, casting a different Rite, negotiating).
    \item \textbf{Tick 2:} The spirit is forced to cross a [WARD] or sacred threshold.
    \item \textbf{Tick 1:} The spirit is wounded (any Harm).
    \item \textbf{Tick 1 per round:} The spirit acts against its nature (fire spirit underwater, sea spirit in a desert).
\end{itemize}

\section{Making Spirits Feel Like Characters}
Before the session, write one sentence for each likely spirit:
\begin{itemize}
    \item “Emberling wants to burn something pretty. If you let it, it will be happy.”
    \item “Drowned Helmsman wants to feel the tiller again. Let it guide your ship, and it will cooperate.”
    \item “Stone-Father wants to hear a new secret. Tell it something it does not know, and it may extend its service.”
\end{itemize}
If the summoner engages with the spirit’s want, reduce Leash ticks or grant a Boon. If ignored, ticks may accelerate.

\section{Summoning as Story Hook}
Every summoning risks something escaping. Use the spirit’s Departure (Leash fills) as a story thread:
\begin{itemize}
    \item The Emberling departs into the chimney. The inn catches fire next week.
    \item The Drowned Helmsman leaves a trail of wet footprints leading to a sea cave where something worse waits.
    \item The Stone-Father sinks into the floor. The building now has a new, angry ghost.
\end{itemize}

\clearpage

\appendix
\chapter{Quick Reference Tables}

\section{Summoning Summary}
\begin{longtable}{p{2.5cm} p{2.5cm} p{8cm}}
\toprule
Culture & Focus & Typical Spirits \\
\midrule
Aeler (Dwarf) & Slate contract & Vein‑spirits, mountain memories \\
Aelaerem (Halfling) & Bread and milk & Brownies, field‑wights \\
Ykrul (Orc) & Drum & Ancestors, totem beasts \\
Linn (Northmen) & Knotted rope & Selkies, drowned crew \\
Zakov (Pirate) & Cursed coin & Drowned mutineers, vengeful dead \\
Valewood & Living wood & Forest shades, beast‑kin \\
Mistlands & Bell & Mist‑wraiths, bell‑echoes \\
\bottomrule
\end{longtable}

\section{Leash Tick Summary}
\begin{itemize}
    \item \textbf{Power use:} +1
    \item \textbf{Summoner splits focus:} +1
    \item \textbf{Crosses [WARD]:} +2
    \item \textbf{Takes Harm:} +1 per level
    \item \textbf{Against nature:} +1 per round
\end{itemize}

\clearpage
\end{document}